\chapter{Potencijalna unapređenja}
\label{ch:unapredjenja}

Prethodna poglavlja zaokružila su tri jezgra i njihov zajednički lanac izvođenja i upotrebe ključa. Ovo poglavlje izlaže nadogradnje koje iz postignutog neposredno slede. Ne ređa ih po jezgrima, nego po tri pravca u kojima bi se sistem dalje razvijao, od puta ka sistemu na ploči, preko pomeranja granice propusnosti puta podataka, do razvoja upravljačke ravni. Iznose se predlozi i njihove očekivane cene, a izmerene vrednosti na koje se predlozi oslanjaju pripadaju poglavlju \ref{ch:rezultati}.

\section{Od simulisanog lanca do sistema na ploči}
\label{sec:unap_sistem}

Lanac iz odeljka \ref{sec:hw_integracija} proveren je u simulaciji, u kojoj ulogu upravljačke logike okvira igra verifikaciono okruženje. Na ploči bi taj posao preuzeo mali namenski automat, i on je prvi korak ka celini. Posao mu je pogodan za hardver, jer su ime funkcije i oznaka namene KMAC okvira fiksirani standardom, a javne delove, so i identitete strana, procesor upiše pri uspostavi sesije, pa je ceo okvir, osim same tajne, poznat pre razmene. Automat bi poznati deo čitao iz male memorije, tajnu $Z$ preuzimao neposredno sa izlaza ECDH jezgra i sastavljeni niz predavao SHA-3 jezgru, čime bi tajni materijal i na ploči ostao unutar programabilne logike, kako predviđa slika \ref{fig:hw_lanac}. Ista razmena pritom može da ponese više od jednog ključa. KMAC iz iste tajne, uz različit kontekst ili veću traženu dužinu izlaza, izvodi onoliko ključeva koliko sistemu treba, po jedan za svaki smer ili za kasnije obnavljanje, pa bi uz automat stajala i mala memorija izvedenih ključeva.
\smallbreak
Sa procesorske strane, upravljački program SHA-3 jezgra iz odeljka \ref{subsec:rez_sha3_drajver} već daje obrazac za ceo lanac, pa bi se proširio komandama za pokretanje razmene i čitanje statusa sesije. Nove veze između jezgara na ploči dobile bi iste skid bafere koji su se na celom putu podataka pokazali bez merljive cene u učestanosti (odeljak \ref{subsec:rez_aes_ceo_put}), uz pun regresioni prolaz posle izmene. Ciljna platforma nudi i delimičnu rekonfiguraciju (engl. \textit{dynamic function exchange}, DFX), a šta ona donosi zavisi od uloge sistema. Na strani klijenta ECDH jezgro radi jednom po sesiji, a u paralelnim konfiguracijama zauzima i preko trećine čipa (odeljak \ref{sec:rez_ecdh}), pa bi se njegova površina posle razmene mogla prepustiti drugom sadržaju, drugom mrežnom kanalu ili nekom drugom algoritmu razmene. Na serverskoj strani, koja razmene vodi za mnogo klijenata istovremeno, jezgro radi neprekidno, pa merodavne postaju latencija skalarnog množenja i broj razmena u sekundi, a izbor se pomera ka najbržim konfiguracijama iz istog odeljka.
\smallbreak
Poglavlje \ref{ch:uvod} predviđa i autentifikaciju učesnika algoritmom ECDSA \cite{fips186}. Njegova realizacija bi merdevine iz odeljka \ref{subsec:hw_ecdh_ladder} upotrebila ponovo za skalarno množenje, ali pored njih traži i aritmetiku po modulu reda tačke, dakle po velikom prostom broju, koje u sadašnjem jezgru nema. Prirodan oblik nadogradnje je zaseban blok za množenje i inverziju po tom modulu, dok bi kriva, tačke i skalarno množenje ostali zajednički. Pun obim je ipak veći, jer ECDSA traži i kvalitetan generator slučajnih brojeva i sabiranje dve nezavisne tačke pri proveri potpisa, pa aritmetički blok pokriva prvi i najveći deo puta. Time bi se zaokružio i bezbednosni model iz poglavlja \ref{ch:osnove}, u kome razmena ključa bez autentifikacije učesnika ne štiti od aktivnog napadača.

\section{Pomeranje granice propusnosti puta podataka}
\label{sec:unap_granica}

Merenja iz odeljka \ref{subsec:rez_aes_ceo_put} pokazuju da put podataka drži blok po taktu i da granicu učestanosti postavlja GCM sloj, ispod koga ceo put podataka staje još za cenu spajanja. Granica pritom ne stoji u šifarskom bloku, jer oba AES bloka sama postižu znatno više od celog puta (odeljak \ref{subsec:rez_aes_resursi}). Stoji u povratnoj petlji GHASH akumulatora, koju registar ne može da preseče a da granica u bitima po taktu ne padne, jer svaki blok zavisi od prethodnog. Rad je zato petlju ipak presekao tamo gde je cilj propusnosti niži, jer dvotaktna organizacija skraćuje kritičnu putanju i olakšava sklapanje sistema, po ceni prepolovljene granice (odeljak \ref{subsec:rez_aes_ceo_put}). Naredni oblik akumulacije istu ideju zadržava bez te cene.
\smallbreak
Povratna petlja ipak može da se ubrza, ali tek zajedno sa oblikom akumulacije. Nad parom blokova akumulacija uz potključ $H$ koristi i $H^2$, pa množenje sa $H$ izlazi iz povratne petlje, a ona se zatvara kroz množač sa $H^2$ jednom u dva takta. U nju zato staje množač sa množenjem u dva takta, onaj koji u GCM sloju sam postiže 302,18 MHz (tabela \ref{tab:ghash_impl}), a koji je do sada polovio propusnost baš zato što je petlju zatvarao svaki takt. Dva množača sa prepolovljenim putanjama tako zadržavaju blok po taktu, a oba šifarska bloka imaju prostora da prate, protočni sa 401,12 MHz, a višejezgarni sa 301,93. Ceo put podataka sa dvotaktnim množačem već dostiže oko 260 MHz, pa bi sa ovakvom akumulacijom na istoj učestanosti nosio pun blok po taktu, oko 33 Gb/s umesto sadašnjih 30. Postupak se uopštava na stepene do $H^k$, gde se petlja zatvara jednom u $k$ taktova. Sa više taktova po množenju množač ne mora ni da ostane u čisto paralelnom Karatsubinom obliku iz odeljka \ref{subsec:hw_ghash}, jer tada u petlju staju i digit-serijske i hibridne organizacije, spoj digit-serijskog toka i Karatsubine dekompozicije, koje kraćim kombinacionim putanjama potencijalno postižu još višu učestanost. Ulaz od dva bloka po taktu, uz ista dva množača u jednotaktnom obliku, udvostručio bi i granicu u bitima po taktu, ali tada ide u paru sa udvostručenim šifarskim putem, na primer sa dva protočna niza i raspodelom paketa među njima. Brzine reda 100 Gb/s traže još širu obradu. Samom AES jezgru dovoljna bi bila potpuno paralelna organizacija na magistrali od 256 bitova, ali ceo AES-GCM put, sa autentifikacionim putem koji granicu mora da prati, realno ide na magistralu od 512 bitova, četiri bloka po taktu, kakvom se sto gigabita u sekundi obično i prenosi. Za takve ciljeve ovaj sistem je na sadašnjoj ploči već ambiciozan, pa je predviđen prelazak na noviju i snažniju platformu.
\smallbreak
Kod SHA-3 jezgra model iz odeljka \ref{subsec:rez_sha3_propusnost} pokazuje da pri dve runde po taktu granicu u većini konfiguracija postavlja unos bloka. Šira ulazna magistrala vratila bi granicu na permutaciju, po ceni šireg upisnog puta u ulaznom baferu.

\section{Upravljačka ravan i budućnost razmene ključa}
\label{sec:unap_uprava}

Prirodna nadogradnja upravljačke ravni je provera primljenog javnog ključa. Jezgro već prima obe koordinate tačke, a provera da tačka leži na krivoj traži da se jednačina krive izračuna i uporedi sa nulom: dva množenja, dva kvadriranja, tri sabiranja i poređenje, ukupno mikroprogram od osam operacija po ugledu na završnu konverziju iz odeljka \ref{subsec:hw_ecdh_ladder}, uz jedan statusni bit ka pozivaocu. Cena od približno $2q$ taktova je ispod desetine procenta trajanja skalarnog množenja (odeljak \ref{subsec:rez_ecdh_latencija}), pa bi provera mogla biti stalno uključena. Puna provera po SP 800-56A \cite{sp80056a}, koja potvrđuje i red tačke, košta ceo drugi prolaz merdevina i pripada protokolskom sloju iznad jezgra.
\smallbreak
Postojeća aritmetika može i bolje da se iskoristi. Merdevine i završna konverzija paralelnog jezgra nikada ne rade istovremeno, pa se tri množača iz odeljka \ref{subsec:hw_ecdh_dva} mogu podeliti sa konverzijom, umesto da ona drži sopstveni, četvrti množač. Konverzija bi se pritom i ubrzala, jer se njena množenja daju rasporediti na tri množača, a inverzija može da pozajmi drugi za svoje korake kvadriranja i množenja. Isti mehanizam otvara aritmetiku i ka spolja, jer se jedan množač može izložiti preko postojećeg start i done interfejsa, na primer budućem ECDSA bloku.
\smallbreak
Prostor postoji i u samim algoritmima množenja u polju, najavljenim u odeljku \ref{subsec:hw_ecdh_mul}. Množač GHASH jedinice već pokazuje kako rekurzivna Karatsubina dekompozicija smanjuje broj elementarnih množenja po ceni dublje mreže, a ista zamena važi i za digit-serijski množač polja, gde hibrid klasičnog i Karatsubinog množenja skraćuje kombinacionu putanju jednog digita \cite{kumari_gf2m_mult}. Kraća putanja diže učestanost, a sa njom i prostor da se pri istom taktu širina digita poveća.
\smallbreak
Na duže staze, pomenuto povlačenje binarnih krivih za nove primene (odeljak \ref{subsec:b571}) znači da nadogradnja u nekom trenutku podrazumeva prelazak na krive nad prostim poljima. Aritmetika se time menja iz korena, jer množenje nad prostim poljem nosi prenos, pa umesto XOR mreža dolaze sabirači i namenski DSP blokovi, koje ciljna platforma nudi u velikom broju, a sadašnji dizajn ne koristi nijedan.
\smallbreak
Najdalji korak za ovaj sistem, ali danas već i preporučeni put za razmenu ključa, jeste postkvantna zamena. Standard FIPS 203 \cite{fips203} propisuje ML-KEM kao mehanizam enkapsulacije ključa, pa hibridni tok iz odeljka \ref{sec:hibridni_sistem} zadržava isti oblik, u kome mehanizam razmene daje tajnu, iz tajne se izvodi ključ sesije, a ključ štiti saobraćaj. U takvoj zameni ECDH jezgro ustupa mesto ML-KEM jezgru, dok SHA-3 jezgro dobija na značaju, jer ML-KEM sve svoje unutrašnje funkcije heširanja i proširivanja gradi nad Keccak permutacijom, istom onom koju ovaj rad već realizuje u hardveru.
